\section*{الفصل \ref{ch:b1:cross-sections} --- المقاطع الجيولوجية}

\begin{solution}{exo:b1:cross-sections:1}
عند مقياس 1:25\,000، يمثّل \qty{1}{cm} على الخريطة \qty{250}{m} على الأرض. والمقياس الرأسي هو \qty{1}{cm} لكل \qty{100}{m}، فيكون المقياس الرأسي $250/100 = 2.5$ مرة المقياس الأفقي: $e = 2.5$.
\end{solution}

\begin{solution}{exo:b1:cross-sections:2}
$\gamma = \ang{90} - \ang{60} = \ang{30}$، و$\tan\beta = \tan\ang{40}
\cos\ang{30} = 0.8391\times0.8660 = 0.7267$: $\beta = \ang{36.0}$، فالطبقة تهبط نحو الطرف الشرقي للمقطع. وفي المقطع ذي \omterm{def:b1:spherical-astronomy:horizontal}{السمت}
\ang{000}، يكون $\gamma = \ang{90}$ و$\beta = 0$: فالمقطع يمتد على طول \omterm{def:b1:geological-maps:attitude}{الامتداد} وتظهر الطبقة أفقية.
\end{solution}

\begin{solution}{exo:b1:cross-sections:3}
$t = h\cos\delta = 120\times\cos\ang{12.16} = 120\times0.9776 =
\qty{117.3}{m}$.
\end{solution}

\begin{solution}{exo:b1:cross-sections:4}
$\tan\ang{12} = 0.2126$. مع $e = 2$: $\arctan 0.4251 = \ang{23.0}$؛ ومع
$e = 3$: $\arctan 0.6377 = \ang{32.5}$؛ ومع $e = 5$: $\arctan 1.0628 =
\ang{46.7}$. فالزاوية تنمو أبطأ كثيرًا من $e$ نفسه، ولا تبلغ \ang{90} أبدًا.
\end{solution}

\begin{solution}{exo:b1:cross-sections:5}
(أ) $t = 300\sin\ang{25} = \qty{126.8}{m}$. (ب) $t = 300\sin\ang{35} =
\qty{172.1}{m}$. (ج) $t = 300\sin\ang{15} = \qty{77.6}{m}$. فقد يعني عرض \omterm{def:b1:geological-maps:geological-map}{الانكشاف} نفسه سمكين يختلفان بأكثر من الضعف: ويجب دائمًا أخذ انحدار الأرض في الحسبان.
\end{solution}

\begin{solution}{exo:b1:cross-sections:6}
بوجود \omterm{def:b1:geological-maps:attitude}{اتجاه الانحدار} $\alpha$ \omterm{def:b1:geological-maps:attitude}{والانحدار} $\delta$، تعطي النظرية في المقطع شمالي-جنوبي ($\theta = 0$) $\tan\ang{20} = \tan\delta\cos\alpha$،
وفي المقطع شرقي-غربي ($\theta = \ang{90}$) $\tan\ang{30} =
\tan\delta\cos(\alpha - \ang{90}) = \tan\delta\sin\alpha$. وبالقسمة،
$\tan\alpha = 0.5774/0.3640 = 1.586$، و$\alpha = \ang{057.8}$؛ وبجمع المربعات،
$\tan\delta = \sqrt{0.3640^2 + 0.5774^2} = 0.6825$، و$\delta =
\ang{34.3}$. والانحداران الظاهريان كلاهما أصغر من \omterm{def:b1:geological-maps:attitude}{الانحدار} الحقيقي، كما يجب أن يكون.
\end{solution}

\begin{solution}{exo:b1:cross-sections:7}
$100/\tan\ang{12.16} = 100/0.2155 = \qty{464.1}{m}$. وعلى الخريطة تكون الكنتورات البنيوية للجناح الشرقي أقلّ قليلًا من نصف كيلومتر من بعضها، متباعدة بانتظام ومتوازية: فالجناح مستوٍ.
\end{solution}

\begin{solution}{exo:b1:cross-sections:8}
B1: $524.02 - 488.98 = \qty{35.04}{m}$؛ وB3: $529.98 - 494.93 =
\qty{35.05}{m}$. وفوق B2، قاعدة الصخر الطيني هي قمة الحجر الجيري، التي تعلو قاعدته بمقدار \qty{180}{m}: $280.16 + 180 = \qty{460.16}{m}$،
وهي أعلى من الأرض (\qty{380.64}{m})، ولهذا تآكلت هناك. وبين B1 وB2 يرتفع \omterm{def:b1:geological-maps:geological-map}{التماس} بمقدار $460.16 - 35.04 = \qty{425.12}{m}$
على مسافة \qty{2000}{m}: أي ميل $0.2126$، وانحدار ظاهري قدره \ang{12.0}. أما B1 وB3، على المسافة نفسها من B2 في الجهتين، فتجدان كل \omterm{def:b1:geological-maps:geological-map}{تماس} على الارتفاع نفسه في حدود \qty{1}{cm}: \omterm{def:b1:geological-maps:fold}{فالطية} متماثلة حول مستوٍ محوري شاقولي عند $x = \qty{3.5}{km}$، وجناحاها ينحدران بالتساوي، ولا يقطع صدع المقطع بينهما.
\end{solution}

\begin{solution}{exo:b1:cross-sections:9}
\omterm{def:b1:cross-sections:throw}{الإزاحة الأفقية} $= 150/\tan\ang{50} = 150/1.1918 = \qty{125.9}{m}$. وفي المقطع ذي \omterm{def:b1:spherical-astronomy:horizontal}{السمت} \ang{150}، يكون $\gamma = \ang{30}$ و$\tan\beta = 1.1918\times0.8660 =
1.0321$: يظهر الصدع بانحدار ظاهري قدره \ang{45.9}. والقطعان السطحيان لطبقة أفقية يفصل بينهما \qty{150}{m} شاقوليًا على ذلك الأثر، ومن ثَمّ $150/1.0321 = \qty{145.3}{m}$ على طول المقطع، أكبر من \omterm{def:b1:cross-sections:throw}{الإزاحة الأفقية}، لأن المقطع يعبر الصدع بشكل مائل.
\end{solution}

\begin{solution}{exo:b1:cross-sections:10}
$\cos\gamma = \tan\ang{20}/\tan\ang{40} = 0.3640/0.8391 = 0.4338$، إذن
$\gamma = \ang{64.3}$: المقاطع ذات السموت $\alpha \pm \ang{64.3}$ (وكذلك الاتجاهات المعاكسة، حيث تصعد الطبقة). ويكون \omterm{def:b1:cross-sections:apparent-dip}{الانحدار الظاهري}
\ang{40.0} و\ang{39.6} و\ang{38.3} و\ang{36.0} من أجل $\gamma = 0$، و\ang{10}، و\ang{20}، و\ang{30}، لكنه \ang{22.8} و\ang{16.0} و\ang{8.3} و$0$ من أجل
$\gamma = \ang{60}$، و\ang{70}، و\ang{80}، و\ang{90}. وبالفعل، بتفاضل
$\tan\beta = \tan\delta\cos\gamma$،
$\dfrac{\mathrm d\beta}{\mathrm d\gamma} =
-\dfrac{\tan\delta\sin\gamma}{1 + \tan^2\delta\cos^2\gamma}$، وهو معدوم
عند $\gamma = 0$ ويساوي $-\tan\delta$ عند $\gamma = \ang{90}$: فالمقطع الذي يبتعد قليلًا عن \omterm{def:b1:geological-maps:attitude}{اتجاه الانحدار} يُظهر \omterm{def:b1:geological-maps:attitude}{الانحدار} الحقيقي تقريبًا، أما قرب \omterm{def:b1:geological-maps:attitude}{الامتداد} فكل درجة من \omterm{def:b1:spherical-astronomy:horizontal}{السمت} تغيّر \omterm{def:b1:cross-sections:apparent-dip}{الانحدار الظاهري} بما يقارب درجة.
\end{solution}

\begin{solution}{exo:b1:cross-sections:11}
تنحدر الطبقات نحو الشرق، فالعمودي على التطبّق، متبعًا نحو الأسفل، يميل \ang{12.16} عن الشاقول نحو الغرب، نحو \ang{260.7}.
(أ) البئر الشاقولي يصنع معه الزاوية $\psi = \ang{12.16}$: $176.0/\cos\ang{12.16} = \qty{180.0}{m}$، وهو \omterm{def:b1:cross-sections:thickness}{السمك الشاقولي}. (ب) بئر يغوص بزاوية \ang{60} نحو \ang{260.7} يميل \ang{30} عن الشاقول في الجهة نفسها التي فيها العمودي: $\psi = \ang{30} - \ang{12.16} =
\ang{17.84}$، و$\cos\psi = 0.9519$، وطول اللبّ $176.0/0.9519 =
\qty{184.9}{m}$. (ج) نحو \ang{080.7} يميل على الجهة الأخرى: $\psi = \ang{30} + \ang{12.16} = \ang{42.16}$، و$\cos\psi = 0.7413$، وطول اللبّ $176.0/0.7413 = \qty{237.4}{m}$. وأقصر لبّ، مساويًا $t$، يكون حين يمتدّ البئر على طول العمودي، أي يغوص بزاوية \ang{77.84} نحو \ang{260.7}.
\end{solution}

\begin{solution}{exo:b1:cross-sections:12}
في المستوي الشاقولي شرق-غرب، مع قياس الزوايا من الشرق، تصعد الطبقات العليا نحو الشرق بزاوية \ang{20}، وتهبط الطبقات السفلى نحو الشرق بزاوية \ang{70}. وبتدوير \omterm{def:b1:geological-maps:geological-map}{الانكشاف} حول \omterm{def:b1:geological-maps:attitude}{الامتداد} بزاوية \ang{20}، مع خفض الجانب الشرقي، تصير الطبقات العليا أفقية وتصل الطبقات السفلى إلى $\ang{70} + \ang{20} = \ang{90}$: أي كانت شاقولية حين رُسبت الطبقات العليا. والتاريخ، من الحدث الأقدم: رُسبت الطبقات السفلى مسطحة؛ ثم قُلبت حتى صارت على حافتها، شاقولية، ورُفعت؛ وتآكلت حوافها حتى صارت سطحًا؛ ورُسبت الطبقات العليا عبر ذلك السطح؛ وأُميل الكل بزاوية \ang{20} نحو الغرب؛ ثم قطعت التعرية سطح الأرض اليوم. وهذا تاريخ \omterm{def:b1:geological-maps:geological-map}{انكشاف} هاتون عند سيكار بوينت.
\end{solution}

\begin{solution}{pb:b1:cross-sections:1}
\textbf{1.} السُّمك فوق النفق هو $z - 200$: أكبره عند $x = \qty{1}{km}$، $556 - 200 = \qty{356}{m}$؛ وأصغره عند $x =
\qty{9}{km}$، $339 - 200 = \qty{139}{m}$.

\textbf{2.} \qty{1}{km} من الارتفاع تقيس \qty{4.2}{cm}، و\qty{1}{km} من المسافة \qty{1.4}{cm}: $e = 4.2/1.4 = 3$.

\textbf{3.} $380.64 - 220.48 = \qty{160.16}{m}$.

\textbf{4.} في B2 تقع قاعدة الحجر الجيري عند $380.64 - 100.48 =
\qty{280.16}{m}$، وقاعدة الحجر الرملي عند \qty{160.16}{m}: فالنفق، على \qty{200}{m}، يقع في الحجر الرملي.

\textbf{5.} B1: $524.02 - 668.98 = \qty{-144.96}{m}$؛ وB3: $529.98 - 674.93 =
\qty{-144.95}{m}$.

\textbf{6.} من الغرب إلى الشرق: الطباشير، والصخر الطيني، والحجر الجيري، والحجر الرملي (في لبّ \omterm{def:b1:geological-maps:fold}{الطية المحدبة})، والحجر الجيري، والصخر الطيني، ثم الطباشير مرة أخرى حتى $x = \qty{9}{km}$، عبر \omterm{def:b1:geological-maps:fold}{الطية المقعرة} وعبر الصدع.

\textbf{7.} $\gamma = \ang{80.7} - \ang{90} = \ang{-9.3}$؛ $\tan\beta =
\tan\ang{12.16}\cos\ang{9.3} = 0.2155\times0.9869 = 0.2126$: $\beta =
\ang{12.0}$، نحو الهبوط شرقًا.

\textbf{8.} $(280.16 + 144.95)/2000 = 425.11/2000 = 0.2126$: \ang{12.0}.

\textbf{9.} $\arctan(3\times0.2126) = \arctan 0.6377 = \ang{32.5}$.

\textbf{10.} $\gamma = \ang{120} - \ang{90} = \ang{30}$؛ $\tan\beta =
\tan\ang{60}\cos\ang{30} = 1.732\times0.866 = 1.500$: $\beta = \ang{56.3}$.

\textbf{11.} ينحدر الجناحان نحو \ang{080.7} أو \ang{279.3}، ولا يبعدان إلا \ang{9.3} عن خط المقطع، و$\cos\ang{9.3} = 0.987$؛ وينحدر الصدع نحو \ang{120}، أي على بعد \ang{30} منه، و$\cos\ang{30} = 0.866$.

\textbf{12.} مستوي الصدع هو $z = 250 - \tan\ang{60}\,[(x - 4300)
\sin\ang{120} + (y - 1000)\cos\ang{120}]$. ومع $y = 6000$، يكون $(y - 1000)
\cos\ang{120} = -2500$، إذن $(x - 4300)\sin\ang{120} = 2500 - (z - 250)/
\tan\ang{60}$، و$1/\tan\ang{60} = \tan\ang{30}$.

\textbf{13.} عند $z = 200$: $x = 4300 + (2500 + 50\times0.5774)/0.8660 =
4300 + 2528.9/0.8660 = \qty{7220}{m}$، أي \qty{7.22}{km} من C.

\textbf{14.} \omterm{def:b1:cross-sections:throw}{الإزاحة الأفقية} $= 200/\tan\ang{60} = \qty{115.5}{m}$؛ والفصل على طول C--D $= 200/1.500 = \qty{133.3}{m}$.

\textbf{15.} ينحدر الصدع نحو الشرق-الجنوب الشرقي، وهو عادي: فالجدار العلوي، إلى الشرق، هو الكتلة الهابطة. وعلى مستوى النفق يقع الطباشير على الجانبين: غرب الصدع، طباشير \omterm{def:b1:geological-maps:fold}{الطية المقعرة}؛ وشرقه، طباشير هبطت \qty{200}{m}.

\textbf{16.} يبقى النفق في الطباشير عبر الصدع، فلا يعلن تغيّر الصخر على الوجه شيئًا؛ ولا يُنذر إلا منطقة من الصخر المسحوق، وربما الماء. ويتنبأ المهندسون بموضعه من المقطع (هنا \qty{7.22}{km})، ويحفرون آبارًا تجريبية أمام الآلة كلما اقتربت، ويعدّون دعامات للمنطقة المسحوقة.

\textbf{17.} $668.98 - 488.98 = \qty{180.00}{m}$.

\textbf{18.} $t = 180\cos\ang{12.16} = \qty{176.0}{m}$.

\textbf{19.} على القمة ($x = 3500$) تقع قاعدة الحجر الرملي عند \qty{160.16}{m} (السؤال 3)؛ وتهبط على الجناحين \omterm{def:b1:cross-sections:apparent-dip}{بالانحدار الظاهري} \ang{12.0} من السؤال 7، أي بمقدار $\tan\ang{12.0}$ لكل متر من المسافة الأفقية عن المحور.

\textbf{20.} قاعدة الحجر الجيري تعلو قاعدة الحجر الرملي بمقدار \qty{120}{m}، وتقع عند $280.16 - u\tan\ang{12.0}$ حيث $u = |x - 3500|$، وقمته أعلى بمقدار \qty{180}{m}، عند $460.16 - u\tan\ang{12.0}$. ومع $\tan\ang{12.0} =
0.21256$: تقع القاعدة تحت \qty{200}{m} حين يكون $u > 80.16/0.21256 = 377.1$، وتقع القمة فوقها حين يكون $u < 260.16/0.21256 = 1224.0$.

\textbf{21.} $x = 3500 \pm 377.1$ و$3500 \pm 1224.0$: يكون النفق في الحجر الجيري من \qty{2276}{m} إلى \qty{3123}{m}، ومن \qty{3877}{m} إلى \qty{4724}{m}.

\textbf{22.} على خط أفقي، تبعد قاعدة الحجر الجيري وقمته، المتباعدتان شاقوليًا بمقدار \qty{180}{m} وكلتاهما تنحدر بمقدار $\tan\beta$، مسافة أفقية $h/\tan\beta$: أي $180/0.21256 = \qty{846.8}{m}$ على كل جناح، وهو نفسه، في حدود التقريب، $1224.0 - 377.1 = 846.9$.

\textbf{23.} $2\times846.8 = \qty{1693.7}{m}$: يحفر النفق \textbf{نحو \qty[text-series-to-math=true]{1694}{m} من الحجر الجيري}، أي أكثر من خُمس طوله البالغ \qty{8}{km}، في مقطعين طول كل منهما \qty{847}{m} على جانبي لبّ الحجر الرملي \omterm{def:b1:geological-maps:fold}{للطية} المحدبة.
\end{solution}
